\section{Additional evaluator checks}
\label{app:evaluator-checks}

These checks distinguish failure to obtain a benchmark score from numerical
nonconvergence, and implementation consistency from robustness of a mesh
recipe. They do not change any original grade or optimizer result.

\subsection{Condition-level outcomes after an accepted mesh}

We audited 3{,}845 evaluation requests with accepted meshes and recoverable,
hash-verified raw result documents. Each contains three condition attempts.
This is a request-level audit: repeated requests and reused results remain
included, and the subset is not a new sample of independent solver runs.
It is separate from the aggregate classification of all 6{,}629 requests.

\begin{table}[htbp]
\centering\small
\caption{Failures after mesh acceptance in the audited subset. The final row
counts evaluations with at least one failed condition, not condition attempts.}
\label{tab:flow-audit}
\begin{tabular}{lrrr}
\toprule
Unit & Failed & Total & Rate \\
\midrule
Mach 0.7853 & 742 & 3{,}845 & 19.3\% \\
Mach 0.8195 & 765 & 3{,}845 & 19.9\% \\
Mach 0.8530 & 896 & 3{,}845 & 23.3\% \\
All condition attempts & 2{,}403 & 11{,}535 & 20.8\% \\
Three-condition evaluations & 1{,}226 & 3{,}845 & 31.9\% \\
\bottomrule
\end{tabular}
\end{table}

The 2{,}403 failed condition attempts comprise 1{,}222 timeouts and 1{,}181
execution/infrastructure outcomes. They cannot be described as numerical
nonconvergence alone. Requiring all three conditions to succeed raises the
observed evaluation failure rate to 31.9\%; this is measured directly, without
assuming independence between conditions. The condition rates remain material
and do not establish adequacy of the solver settings or time limits.

\subsection{Source-geometry controls and upstream meshing}

We remeshed 1{,}000 distinct previously successful SuperWing geometries with
the benchmark spacing and the original source spacing. The two cohorts of
500 were separately frozen random selections from eligible source geometries;
the second was selected after the first cohort's results. We therefore give
descriptive combined counts, without a pooled confirmatory interval. We also
tested all 1{,}913 distinct failed geometries recoverable from the benchmark
records with the original, unmodified meshing script released with SuperWing
in the AeroTransformer repository. The script is byte-for-byte identical to
the published version pinned in Appendix~\ref{app:cfdprotocol}, and was run
directly with its original tip-template path, without the WingBench case
preparer or grading orchestration, under matched inputs, MPI ranks, runtime,
and timeout settings.

\begin{table}[htbp]
\centering\small
\caption{Mesh acceptance in source controls and the complete recoverable
failed-geometry census. Acceptance requires
positive minimum cell volume and quality; marginal positive quality is accepted.
These are mesh checks, not completed CFD evaluations.}
\label{tab:mesh-controls}
\begin{tabular}{llrr}
\toprule
Cohort & Mesh setting & Accepted & Tested \\
\midrule
Successful source geometries & Benchmark spacing & 999 & 1{,}000 \\
Same source geometries & Original source spacing & 1{,}000 & 1{,}000 \\
Previously failed geometries & Upstream default (81 layers) & 1 & 1{,}913 \\
Same failed geometries & Diagnostic 129 layers & 605 & 1{,}913 \\
\bottomrule
\end{tabular}
\end{table}

The failed-geometry census is complete: all 30 prespecified batches from the
randomized ordering of 1{,}913 distinct geometries finished by 9 November 2026.
It includes every recoverable geometry identified as mesh-blocked in either
the released request categories or the raw evaluation records. One additional
mesh-failure request has no archived geometry and could not be rerun.
Under the unchanged upstream default, 408 were rejected by mesh gates,
1{,}504 stopped inconclusively, and one produced an accepted mesh.
Thus 1{,}912/1{,}913 remained unsuccessful.
This supports consistency with the published procedure; it does not establish
that the procedure is robust throughout the search box or that failed shapes
are physically infeasible. Rebuilt environments also do not establish exact
historical binary and hardware equivalence.

The paired 129-layer diagnostic is an important qualification: it accepted
605 meshes, including the one accepted by the default, so 604 default failures
became mesh successes. Failures are therefore sensitive to numerical mesh
settings. These meshes have not been used to replace benchmark grades, and
mesh acceptance alone says nothing about subsequent convergence or $L/D$.
Improving the recipe requires separate flow validation and a new protocol
version. The current evidence narrows an implementation-mismatch explanation
but does not establish a universal ground-truth feasibility evaluator.

The portable tables \texttt{accepted\_mesh\_flow\_audit.csv} and
\texttt{mesh\_validation\_checks.csv}, with source fingerprints in
\texttt{revision\_evidence\_provenance.json}, preserve the counting units and
cohort membership. \texttt{analysis/21\_revision\_evidence.py} reproduces the
tables above without new CFD runs.

\subsection{Surface appearance and automatic volume meshing}
\label{app:automatic-meshing}

A wing surface and a valid surrounding volume mesh are different objects.
The published SuperWing script constructs surface blocks and uses pyHyp to
march them toward the far field under fixed settings. Local curvature can
make marching lines approach or cross; pyHyp documents smoothing controls
for preventing such crossings in concave regions \citep{mdolabpyhyp}.
This is a numerical limitation that can arise with a correctly implemented
meshing method. A different topology, tuned settings or manual mesh construction
could succeed on a surface for which the configured automatic recipe fails.
Such alternatives must be validated and specified as a separate protocol;
the benchmark does not perform unrecorded, wing-specific manual repairs.

